% ========== DEPENDENCY MANAGEMENT ==========

\section{Dependency Management}
\label{sec:dependency-management}

\lettrine{D}{ependency} management in Symphony's extension ecosystem requires sophisticated algorithms and policies to handle the complex web of relationships between extensions, system components, and external dependencies. The system must balance compatibility, security, performance, and user experience while enabling the rich extension interactions that make Symphony's AI-first approach possible.

\subsection{Semantic Versioning Framework}
\label{subsec:semantic-versioning-framework}

Symphony adopts and extends semantic versioning (SemVer 2.0) with additional metadata and constraints specific to extension ecosystems and AI model compatibility.

\subsubsection{Extended SemVer Format}
\label{subsubsec:extended-semver-format}

The Symphony versioning system uses the format: \texttt{MAJOR.MINOR.PATCH[-PRERELEASE][+BUILD]}

\begin{table}[h]
\centering
\caption{Symphony Version Component Semantics}
\label{tab:version-semantics}
\begin{tabular}{@{}lll@{}}
\toprule
\textbf{Component} & \textbf{Increment Trigger} & \textbf{Compatibility Impact} \\
\midrule
MAJOR & Breaking API changes & Incompatible with previous versions \\
MINOR & New features, backward compatible & Compatible, new capabilities \\
PATCH & Bug fixes, security patches & Compatible, no new features \\
PRERELEASE & Alpha, beta, release candidate & Testing and validation phases \\
BUILD & Build metadata, commit hash & No semantic meaning \\
\bottomrule
\end{tabular}
\end{table}

\subsubsection{AI Model Versioning}
\label{subsubsec:ai-model-versioning}

Instrument extensions (AI models) require special versioning considerations due to their non-deterministic nature and evolving capabilities:

\begin{lstlisting}[language=toml,caption=AI Model Version Declaration]
[package]
version = "2.1.3-beta.1+gpt4.20240315"

[model.versioning]
# Model-specific version information
model_version = "gpt-4-turbo-2024-04-09"
training_cutoff = "2024-04-01"
capability_version = "2.1"  # Semantic capability versioning

# Compatibility declarations
api_compatibility = "^2.0.0"  # Compatible API versions
behavior_stability = "stable"  # stable | evolving | experimental
determinism_level = "low"      # high | medium | low
\end{lstlisting}

\subsection{Dependency Resolution Algorithm}
\label{subsec:dependency-resolution-algorithm}

Symphony's dependency resolution algorithm combines constraint satisfaction with heuristic optimization to find optimal extension configurations while respecting user preferences and system constraints.

\subsubsection{Resolution Strategy}
\label{subsubsec:resolution-strategy}

The resolution process follows a multi-phase approach that balances correctness with performance:

\begin{figure}[htbp]
\centering
\begin{tikzpicture}[
    phase/.style={rectangle, rounded corners, draw=brandPrimary, fill=brandPrimary!10, text width=3cm, text centered, minimum height=1.2cm},
    decision/.style={diamond, draw=brandSecondary, fill=brandSecondary!10, text width=2.5cm, text centered, minimum height=1cm},
    flow/.style={->, >=stealth, thick, brandAccent}
]

% Resolution phases
\node[phase] (collect) {Collect Requirements};
\node[phase, below=1.5cm of collect] (analyze) {Analyze Constraints};
\node[decision, below=1.5cm of analyze] (solvable) {Solvable?};
\node[phase, left=2.5cm of solvable] (negotiate) {Negotiate Conflicts};
\node[phase, below=1.5cm of solvable] (optimize) {Optimize Solution};
\node[phase, below=1.5cm of optimize] (validate) {Validate Result};

% Flow arrows
\draw[flow] (collect) -- (analyze);
\draw[flow] (analyze) -- (solvable);
\draw[flow] (solvable) -- node[above] {No} (negotiate);
\draw[flow] (negotiate) |- (analyze);
\draw[flow] (solvable) -- node[right] {Yes} (optimize);
\draw[flow] (optimize) -- (validate);

% Side annotations
\node[right=1cm of collect, text width=3cm] {\small Parse manifests, extract dependencies};
\node[right=1cm of analyze, text width=3cm] {\small Build constraint graph, detect conflicts};
\node[right=1cm of optimize, text width=3cm] {\small Select optimal versions, minimize conflicts};

\end{tikzpicture}
\caption{Dependency Resolution Algorithm Phases}
\label{fig:dependency-resolution-phases}
\end{figure}

\subsubsection{Constraint Satisfaction}
\label{subsubsec:constraint-satisfaction}

The algorithm models dependency resolution as a constraint satisfaction problem (CSP) with the following components:

\begin{compactlist}
\item \textbf{Variables}: Extension versions to be selected
\item \textbf{Domains}: Available versions for each extension
\item \textbf{Constraints}: Version compatibility requirements
\item \textbf{Objective}: Minimize conflicts while maximizing user preferences
\end{compactlist}

\begin{lstlisting}[language=rust,caption=Dependency Resolution Core Algorithm]
pub struct DependencyResolver {
    registry: ExtensionRegistry,
    constraints: ConstraintGraph,
    preferences: UserPreferences,
}

impl DependencyResolver {
    pub fn resolve(&self, requirements: &[Requirement]) -> ResolutionResult {
        // Phase 1: Collect all transitive dependencies
        let all_deps = self.collect_transitive_dependencies(requirements)?;
        
        // Phase 2: Build constraint satisfaction problem
        let csp = self.build_constraint_problem(&all_deps)?;
        
        // Phase 3: Solve with backtracking and heuristics
        let solution = self.solve_csp(&csp)?;
        
        // Phase 4: Optimize solution for user preferences
        let optimized = self.optimize_solution(solution, &self.preferences)?;
        
        // Phase 5: Validate final solution
        self.validate_solution(&optimized)
    }
    
    fn solve_csp(&self, csp: &ConstraintProblem) -> Option<Solution> {
        // Use arc consistency + backtracking with intelligent heuristics
        let mut solver = CSPSolver::new(csp);
        solver.set_variable_ordering(MostConstrainedFirst);
        solver.set_value_ordering(PreferStableVersions);
        solver.solve()
    }
}
\end{lstlisting}

\subsection{Conflict Resolution}
\label{subsec:conflict-resolution}

When dependency conflicts arise, Symphony employs multiple resolution strategies that prioritize system stability while respecting user intentions.

\subsubsection{Conflict Types}
\label{subsubsec:conflict-types}

\begin{table}[h]
\centering
\caption{Dependency Conflict Types and Resolution Strategies}
\label{tab:conflict-types}
\begin{tabular}{@{}lll@{}}
\toprule
\textbf{Conflict Type} & \textbf{Description} & \textbf{Resolution Strategy} \\
\midrule
Version Incompatibility & Incompatible version ranges & Find compatible intersection \\
Mutual Exclusion & Extensions cannot coexist & Choose based on priority \\
Resource Competition & Competing for same resources & Implement sharing or queuing \\
API Conflicts & Conflicting API requirements & Mediate through adapters \\
Circular Dependencies & Circular dependency chains & Break cycles with lazy loading \\
\bottomrule
\end{tabular}
\end{table}

\subsubsection{Resolution Strategies}
\label{subsubsec:resolution-strategies}

Symphony implements multiple conflict resolution strategies that can be applied individually or in combination:

\begin{infobox}[title=Conflict Resolution Hierarchy]
\textbf{1. Automatic Resolution:} Use algorithm to find compatible versions \\
\textbf{2. User Mediation:} Present options and let user choose \\
\textbf{3. Graceful Degradation:} Disable conflicting features \\
\textbf{4. Alternative Suggestions:} Recommend compatible alternatives
\end{infobox}

\begin{expandedlist}
\item \textbf{Version Negotiation}: Automatically find version ranges that satisfy all constraints
\item \textbf{Priority-Based Selection}: Choose extensions based on user-defined or system-calculated priorities
\item \textbf{Feature Isolation}: Isolate conflicting features while maintaining core functionality
\item \textbf{Adapter Injection}: Insert compatibility adapters to mediate between conflicting APIs
\item \textbf{Lazy Resolution}: Defer conflict resolution until runtime when more context is available
\end{expandedlist}

\subsubsection{User-Guided Resolution}
\label{subsubsec:user-guided-resolution}

When automatic resolution fails, Symphony presents clear options to users:

\begin{figure}[htbp]
\centering
\begin{tikzpicture}[
    option/.style={rectangle, rounded corners, draw=brandSecondary, fill=brandSecondary!10, text width=4cm, text centered, minimum height=1.5cm},
    impact/.style={ellipse, draw=brandAccent, fill=brandAccent!10, text width=2.5cm, text centered, minimum height=0.8cm}
]

% Resolution options
\node[option] (keep) {Keep Current Versions \\ \small Maintain stability};
\node[option, right=1.5cm of keep] (upgrade) {Upgrade Dependencies \\ \small Get new features};
\node[option, below=2cm of keep] (alternative) {Use Alternative Extension \\ \small Different approach};
\node[option, below=2cm of upgrade] (disable) {Disable Conflicting Features \\ \small Reduced functionality};

% Impact indicators
\node[impact, above=0.5cm of keep] {\small Low Risk};
\node[impact, above=0.5cm of upgrade] {\small Medium Risk};
\node[impact, below=0.5cm of alternative] {\small Learning Curve};
\node[impact, below=0.5cm of disable] {\small Feature Loss};

\end{tikzpicture}
\caption{User Conflict Resolution Options}
\label{fig:user-conflict-resolution}
\end{figure}

\subsection{Transitive Dependencies}
\label{subsec:transitive-dependencies}

Managing transitive dependencies requires careful analysis of the entire dependency graph to ensure consistency and avoid unexpected conflicts.

\subsubsection{Dependency Graph Analysis}
\label{subsubsec:dependency-graph-analysis}

Symphony builds and analyzes the complete dependency graph to understand relationships and potential issues:

\begin{table}[h]
\centering
\caption{Dependency Graph Metrics}
\label{tab:dependency-graph-metrics}
\begin{tabular}{@{}lll@{}}
\toprule
\textbf{Metric} & \textbf{Purpose} & \textbf{Threshold} \\
\midrule
Graph Depth & Detect deep dependency chains & Warning at depth > 10 \\
Cycle Detection & Identify circular dependencies & Error on any cycle \\
Fan-out & Extensions with many dependencies & Warning at fan-out > 20 \\
Shared Dependencies & Common dependency usage & Optimize for sharing \\
Version Spread & Range of versions for same extension & Minimize spread \\
\bottomrule
\end{tabular}
\end{table}

\subsubsection{Transitive Conflict Detection}
\label{subsubsec:transitive-conflict-detection}

The system proactively detects potential conflicts in transitive dependencies:

\begin{successbox}
\textbf{Early Detection:} Symphony analyzes transitive dependencies during manifest validation, catching potential conflicts before installation rather than at runtime.
\end{successbox}

\begin{compactlist}
\item \textbf{Version Range Analysis}: Check if transitive version requirements can be satisfied
\item \textbf{Capability Conflicts}: Detect when transitive dependencies have conflicting capabilities
\item \textbf{Resource Conflicts}: Identify potential resource competition in dependency chains
\item \textbf{Security Implications}: Analyze security implications of transitive trust relationships
\end{compactlist}

\subsection{Version Constraints}
\label{subsec:version-constraints}

Symphony supports rich version constraint expressions that enable precise dependency specification while maintaining flexibility for ecosystem evolution.

\subsubsection{Constraint Operators}
\label{subsubsec:constraint-operators}

\begin{table}[h]
\centering
\caption{Version Constraint Operators}
\label{tab:version-constraint-operators}
\begin{tabular}{@{}lll@{}}
\toprule
\textbf{Operator} & \textbf{Meaning} & \textbf{Example} \\
\midrule
\texttt{=} & Exact version & \texttt{=1.2.3} \\
\texttt{>=} & Minimum version & \texttt{>=1.0.0} \\
\texttt{<=} & Maximum version & \texttt{<=2.0.0} \\
\texttt{\textasciicircum} & Compatible within major & \texttt{\textasciicircum 1.2.0} \\
\texttt{\textasciitilde} & Compatible within minor & \texttt{\textasciitilde 1.2.3} \\
\texttt{*} & Wildcard matching & \texttt{1.2.*} \\
\texttt{||} & Logical OR & \texttt{>=1.0.0 || >=2.0.0} \\
\texttt{\&\&} & Logical AND & \texttt{>=1.0.0 \&\& <2.0.0} \\
\bottomrule
\end{tabular}
\end{table}

\subsubsection{Advanced Constraint Features}
\label{subsubsec:advanced-constraint-features}

Symphony extends basic version constraints with additional features for complex scenarios:

\begin{lstlisting}[language=toml,caption=Advanced Version Constraints]
[dependencies]
# Basic constraints
rust_analyzer = "^1.0.0"

# Complex constraints with conditions
ai_model = { version = ">=2.0.0", features = ["streaming", "function-calls"] }

# Platform-specific constraints
[dependencies.windows]
windows_integration = "^1.5.0"

[dependencies.macos]
macos_integration = "^1.3.0"

# Conditional dependencies
[dependencies.optional]
gpu_acceleration = { version = "^1.0.0", optional = true }
enterprise_features = { version = "^2.0.0", requires = "enterprise_license" }

# Conflict declarations
[conflicts]
legacy_rust_support = "<1.0.0"  # Cannot coexist with old versions
alternative_ai_model = "*"      # Mutually exclusive
\end{lstlisting}

\subsubsection{Constraint Validation}
\label{subsubsec:constraint-validation}

The system validates version constraints at multiple levels:

\begin{alertbox}
\textbf{Constraint Validation:} Invalid or unsatisfiable constraints are detected during manifest parsing, preventing extensions with impossible dependencies from entering the ecosystem.
\end{alertbox}

\begin{compactlist}
\item \textbf{Syntax Validation}: Ensure constraint expressions are syntactically correct
\item \textbf{Semantic Validation}: Verify constraints are logically consistent
\item \textbf{Satisfiability Check}: Confirm constraints can be satisfied with available versions
\item \textbf{Performance Impact}: Analyze constraint complexity for resolution performance
\end{compactlist}

The dependency management system provides the foundation for Symphony's rich extension ecosystem, enabling complex interactions while maintaining system stability and user control over their development environment.